\begin{table}[t]
\centering
\footnotesize
\setlength{\tabcolsep}{6pt}
\begin{threeparttable}
\caption{%
\textbf{Role of supervision signal.} Pointing to bounding box centers (Gaussian Centroid) during training is inferior to the segmentation objective across all downstream evaluations. Segmentation provides \model{} with better semantic awareness and tighter multimodal coupling, as evidenced by the much improved fine-grained language understanding on PODS. 
}
\label{tab:supervision}
\begin{tabular}{@{} l
                S[table-format=2.1]
                S[table-format=2.1]
                S[table-format=2.1]
                S[table-format=2.1]@{}}
\toprule
{Supervision Signal}
  & {FG-CLS $\uparrow$}
  & {ADE20k $\uparrow$}
  & {CORE\ $\uparrow$}
  & {PODS $\uparrow$} \\
\midrule
Gaussian Centroid & {61.8} & {49.8} & {94.3} & {37.6} \\
\rowcolor{\methodcolor}
Segmentation Mask & \textbf{76.2}  & \textbf{{55.1}} & \textbf{94.9} & \textbf{50.6} \\
\rowcolor{\methodcolor}
% 5 & RoBERTa & \textbf{Large (ours)} & \textbf{77.0} & \textbf{93.4} & \textbf{55.3} & \textbf{55.3} \\
\bottomrule
\end{tabular}
% \begin{tablenotes}[flushleft]\footnotesize
% \scriptsize
% \item ViT scaling w/ RoBERTa-Large; RoBERTa scaling w/ DINOv2 ViT-B.
% % All models trained for 500k iterations on the full data mixture.
% % PODS uses descriptive text prompts.
% \end{tablenotes}
\end{threeparttable}
\end{table}
